\documentclass[11pt]{article}
\usepackage[T1]{fontenc}\usepackage[utf8]{inputenc}\usepackage{lmodern}\usepackage[margin=2.5cm]{geometry}\usepackage{microtype}\usepackage{hyperref}
% generated by code/experiments/aggregate.py -- do not edit
\newcommand{\ATwoK}{7}
\newcommand{\ATwoN}{7}
\newcommand{\BudgetLsOne}{27}
\newcommand{\BudgetLsTen}{36}
\newcommand{\BudgetLsTenth}{19}
\newcommand{\BudgetMainLsTenth}{9}
\newcommand{\BudgetMainNlTenth}{4}
\newcommand{\BudgetNlOne}{28}
\newcommand{\BudgetNlTen}{36}
\newcommand{\BudgetNlTenth}{10}
\newcommand{\ClsLsAtFivetwelve}{12}
\newcommand{\ClsLsAtOnetwentyeight}{7}
\newcommand{\ClsLsAtSixtyfour}{8}
\newcommand{\ClsLsAtTentwentyfour}{13}
\newcommand{\ClsLsAtThirtytwo}{5}
\newcommand{\ClsLsAtTwofiftysix}{9}
\newcommand{\ClsNlAtFivetwelve}{5}
\newcommand{\ClsNlAtOnetwentyeight}{8}
\newcommand{\ClsNlAtSixtyfour}{5}
\newcommand{\ClsNlAtTentwentyfour}{8}
\newcommand{\ClsNlAtThirtytwo}{5}
\newcommand{\ClsNlAtTwofiftysix}{4}
\newcommand{\ClsTieAtFivetwelve}{7}
\newcommand{\ClsTieAtOnetwentyeight}{9}
\newcommand{\ClsTieAtSixtyfour}{11}
\newcommand{\ClsTieAtTentwentyfour}{3}
\newcommand{\ClsTieAtThirtytwo}{14}
\newcommand{\ClsTieAtTwofiftysix}{11}
\newcommand{\CpuDirectHi}{139}
\newcommand{\CpuDirectLo}{116}
\newcommand{\CpuGp}{44}
\newcommand{\CpuPodmlp}{86}
\newcommand{\DrawAgree}{83\%}
\newcommand{\DrawSign}{99\%}
\newcommand{\ExpBChanged}{15}
\newcommand{\ExpBImproved}{90\%}
\newcommand{\ExpBRetention}{88\%}
\newcommand{\ExpBShift}{0.24}
\newcommand{\ExpBTotal}{48}
\newcommand{\ExpFGain}{0.47}
\newcommand{\ExpFKept}{5}
\newcommand{\ExpFN}{7}
\newcommand{\ExtMedianExt}{-0.03}
\newcommand{\ExtMedianInterp}{0.16}
\newcommand{\FomMedianMs}{6.0}
\newcommand{\FomN}{36}
\newcommand{\FullCellsN}{28}
\newcommand{\FullOverPod}{0.95}
\newcommand{\HfourAB}{10}
\newcommand{\HfourAC}{1}
\newcommand{\HfourAP}{0.006}
\newcommand{\HoneAMedian}{-0.059}
\newcommand{\HoneAP}{0.874}
\newcommand{\HoneBK}{5}
\newcommand{\HtwoPHolm}{0.145}
\newcommand{\IsoOverArd}{1.92}
\newcommand{\KrrNarrowN}{28}
\newcommand{\KrrNarrowRatio}{13}
\newcommand{\KrrNarrowShareThree}{18\%}
\newcommand{\KrrWideEightBeat}{61\%}
\newcommand{\KrrWideEightMed}{6.5}
\newcommand{\KrrWideEightShare}{21\%}
\newcommand{\KrrWideThirtytwoBeat}{64\%}
\newcommand{\KrrWideThirtytwoMed}{2.7}
\newcommand{\KrrWideThirtytwoShare}{57\%}
\newcommand{\LeakErrorRatio}{229}
\newcommand{\LeakMaxMatches}{0}
\newcommand{\LeakPlantedFlagged}{20}
\newcommand{\MetricAgreeMed}{95\%}
\newcommand{\NMainProblems}{24}
\newcommand{\NOwnProblems}{36}
\newcommand{\NlClearTieBandZeroonefive}{6}
\newcommand{\NlClearTieBandZerothree}{3}
\newcommand{\NlClearTieBandZerozerofive}{10}
\newcommand{\NnFloorF}{$2.0\times10^{-4}$}
\newcommand{\PLeFourLs}{20}
\newcommand{\PLeFourN}{30}
\newcommand{\PLeFourNl}{8}
\newcommand{\PoneP}{0.18}
\newcommand{\PoneRho}{-0.38}
\newcommand{\PooledRmseBase}{1.07}
\newcommand{\PooledRmseCv}{0.62}
\newcommand{\PtwoGain}{0.60}
\newcommand{\PtwoPHolm}{0.051}
\newcommand{\SeedAllLs}{80\%}
\newcommand{\SeedAllNl}{57\%}
\newcommand{\SoneCorrect}{6}
\newcommand{\SoneN}{10}

\setlength{\parindent}{0pt}\setlength{\parskip}{0.7em}
\begin{document}
[Date]

To the Editor of [journal name --- to be chosen by the authors],

\textbf{Submission: ``Choosing a Fast Stand-In for a Slow Simulation: a Selection Tool and a 44-Problem Study''}

\textbf{The practical problem.} Engineers who rely on a slow simulation can afford only a limited number of runs, and then have to decide which fast stand-in model to train and whether more runs would be worth the cost. Published comparisons of Gaussian-process, reduced-basis and neural-network stand-ins disagree, are often run under unequal tuning, and do not give a procedure a practitioner can run on their own data.

\textbf{What the paper offers.} (1) An open tool, SurroSelect, that chooses a stand-in by cross-validation, reports which models are tied, checks whether the POD basis limits accuracy, and forecasts the error for more runs. (2) A pre-specified evaluation of the tool (criteria committed to version control before the evaluation runs) on \NTpProblems{} problems solved by an independent finite-element library, against a tuned neural network, and on two real measured datasets. (3) A study of \NAllProblems{} problems describing where each family wins, with controls for the reduced basis, the regressor and the width of the hyper-parameter search.

\textbf{The evidence, in numbers.} The criteria were first tested on \NTpProblems{} problems; the tool was then frozen, with predictions written down, and tested on \NTpProblems{} further ones. On the further problems selection regret against the oracle had median \SvSoneMedian{} (95\% interval \SvSoneLo--\SvSoneHi), and the tuned network was better than the tool by more than 25\% in \SvSfourK{} of \SvSfourN{} cells. Several results went against the tool or against our predictions: it did not match the benchmark's own selection among methods that include networks (in either batch), the learning-curve forecast was within a factor two in only \SvpBoneShare{} of \SvpBoneN{} cases over all sixteen problems, the basis check raised no flag on the further problems, and the benchmark finding that a reduced-basis model wins at large budgets held on \SvPtK{} of \SvPtN{} of them where we had predicted that it would not. Adding a tuned network to the pool helped on one of two wave problems, at several times the cost, and on a solver taking \SvETsim{} seconds per run the surrogate was cheaper than the solver after about \SvEBreakTwofiftysix{} predictions. Limitations are listed in the paper: the problems are simulated and their solvers are fast, the real-data check uses scalar noisy outputs, the networks were tuned with a limited budget, and the forecast is a rough extrapolation.

\textbf{Why this journal.} [To be written by the authors after choosing the journal and checking its scope, article types, indexing and author guidelines; the project record does not contain a verified match for the authors' constraints.]

\textbf{Statements.} The manuscript is original and is not under review elsewhere [authors to confirm]. All data are generated by the code in the project repository or are public, openly licensed measurements (CC BY 4.0); the tool is released under the MIT licence; a permanent archive with a DOI [to be provided]. The study design, code, analyses and drafting used an AI system (Claude) under the authors' direction, as disclosed in the acknowledgments; the internal reviews were AI-agent reviews, not human peer review; all authors must approve the submission.

Sincerely,\\\relax
[Author names, affiliations, corresponding-author e-mail --- to be supplied by the authors]
\end{document}
